% ID: SEL-D0032-OB-E01
% ESTADO: validada. S033, cotejada con examen y criterios.
Sean $x$ las empresas e $y$ los particulares. Deben cumplirse
$x\ge25$, $y\ge2x$ y $x+y\le120$. El beneficio es $386x+229y$.
Las fronteras se cortan en $(25,50)$, $(25,95)$ y $(40,80)$:
en $x=25$ se obtienen los dos primeros y $2x=120-x$
da el tercero.
\begin{center}
\setlength{\unitlength}{5mm}
\begin{picture}(14,13)(-1,-1)
\qbezier(0,0)(5,0)(10,0)\put(10.2,0){$x$}
\qbezier(0,0)(0,5)(0,10)\put(0,10.2){$y$}
\qbezier(5,5)(5,7.25)(5,9.5)
\qbezier(5,9.5)(6.5,8.75)(8,8)
\qbezier(8,8)(6.5,6.5)(5,5)
\put(5,4.7){\makebox(0,0)[t]{$(25,50)$}}
\put(5,9.7){\makebox(0,0)[b]{$(25,95)$}}
\put(8.2,8){$(40,80)$}
\end{picture}
\end{center}
Los beneficios en esos vértices son $21100$, $31405$ y $33760$
euros, respectivamente. Una función lineal alcanza el máximo
en algún vértice. El mejor tiene coordenadas enteras, así que
también es óptimo si sólo se cuentan clientes completos.
La combinación es \textbf{40 empresas y 80 particulares},
con $\boxed{33760}$ euros de beneficio.
